\documentclass[11pt]{article}
\usepackage[T1]{fontenc}
\usepackage[utf8]{inputenc}
\usepackage{amsmath,amssymb,amsthm}
\usepackage{graphicx}
\usepackage[margin=1in]{geometry}
\usepackage[colorlinks=true,linkcolor=blue,citecolor=blue,urlcolor=blue]{hyperref}
\emergencystretch=3em

\newtheorem{theorem}{Theorem}
\newtheorem{lemma}{Lemma}
\newtheorem{proposition}{Proposition}
\newtheorem{remark}{Remark}

\newcommand{\hhilb}{\mathfrak h}
\newcommand{\HH}{\mathcal H}
\newcommand{\FF}{\mathcal F}
\newcommand{\G}{\mathfrak H}
\newcommand{\dom}{\operatorname{Dom}}
\newcommand{\ran}{\operatorname{Ran}}
\newcommand{\kerop}{\operatorname{Ker}}
\newcommand{\supp}{\operatorname{supp}}
\newcommand{\dd}{\,d}

\title{A Spectral-Calculus Extension of the Mild Generalized Spin--Boson Theorem}
\author{run fa\_banach\_001, lane 3}
\date{2026-07-19}

\begin{document}
\maketitle

\paragraph{Status.}
\textbf{partial\_result\_likely\_valid.} This packet proves a conservative
extension of the source paper's Theorem 1.3 to a separable infinite-dimensional
spin space with bounded commuting normal interaction operators. It does not
address the source paper's Case 2 or Case 3 ultraviolet-divergence problems.

\paragraph{Source signal.}
The source paper is Lill--Lonigro \cite{LL}, arXiv:2307.14727. On page 6,
Remark 1.5 says that the finite-dimensional proof should extend to
infinite-dimensional spin space when the interaction operators are bounded and
satisfy Assumption 1.1, but that Proposition 4.4 has to be replaced by a
spectral-calculus argument.

\begin{center}
\includegraphics[width=0.96\linewidth]{figures/open_problem_crop.png}
\end{center}

\paragraph{Problem Extracted.}
Give the omitted spectral-calculus proof of the infinite-dimensional spin-space
analogue of Theorem 1.3 in the mild-divergence Case 1, at least under the same
commuting-normal/common-kernel hypotheses on the interaction operators.

\section*{The Result}

Let $\HH$ be the one-particle Hilbert space, let $\omega\ge m_0>0$ be the
dispersion relation, and let $\FF=\FF^+(\HH)$ be the bosonic Fock space. For
$r\in\mathbb R$ write $\HH_r$ and $\FF_r$ for the Hilbert scales used in
\cite{LL}. Let $\hhilb$ be a separable Hilbert space and let
$K\in\mathcal B(\hhilb)$ be self-adjoint. After adding a scalar constant, assume
for notational simplicity that
\[
H_{\rm free}=K\otimes I_{\FF}+I_{\hhilb}\otimes d\Gamma(\omega)
\]
is nonnegative. This harmless shift does not change the domain formula below.
Put $\G=\hhilb\otimes\FF$ and let $\G_r$ be the Hilbert scale generated by
$H_{\rm free}$.

\begin{theorem}
Fix $0<s\le 1$. Let $f_1,\ldots,f_N\in \HH_{-s}\setminus\HH$ be
$\HH$-independent in the sense of Definition 1.2 of \cite{LL}. Let
$B_1,\ldots,B_N\in\mathcal B(\hhilb)$ be bounded normal operators such that
$B_jB_\ell=B_\ell B_j$ for all $j,\ell$ and
\[
\bigcap_{j=1}^N \kerop B_j=\{0\}.
\]
Define
\[
A=\sum_{j=1}^N B_j^*\otimes a(f_j),\qquad
A^*=\sum_{j=1}^N B_j\otimes a^*(f_j),
\]
where the singular creation and annihilation operators are interpreted as in
\cite[Proposition 4.1]{LL}. Then the formal Hamiltonian
\[
H=H_{\rm free}+A^*+A
\]
has the same conclusions as \cite[Theorem 1.3]{LL}:
\[
\dom(H)=
\left\{\Psi\in\G_{2-s}:
\left[1+(H_{\rm free}-z_0)^{-1}A^*\right]\Psi\in\dom(H_{\rm free})\right\},
\]
for any real $z_0\in\rho(H_{\rm free})$, and $H$ is self-adjoint on this domain.
Moreover, for each $j$ there are $f_{j,n}\in\HH$ with
$\|f_{j,n}-f_j\|_{\HH_{-s}}\to0$ such that, with
\[
A_n=\sum_{j=1}^N B_j^*\otimes a(f_{j,n}),\qquad
A_n^*=\sum_{j=1}^N B_j\otimes a^*(f_{j,n}),
\]
the cutoff Hamiltonians $H_{\rm free}+A_n^*+A_n$ converge to $H$ in norm
resolvent sense.
\end{theorem}

\section*{Proof Intuition}

The finite-dimensional source proof uses the simultaneous diagonalization of
the matrices $B_j$. On each common eigenspace, $A$ reduces to one annihilation
operator $a(\sum b_jf_j)$; the source's Lemma 4.2 then gives dense range,
because $\HH$-independence keeps this linear combination outside $\HH$ whenever
the coefficient vector is nonzero.

For bounded commuting normal operators, the joint spectral theorem gives the
same picture with direct integrals instead of a finite direct sum. The common
kernel condition says that the joint spectral point $(0,\ldots,0)$ has zero
spectral projection. Thus almost every fiber sees a nonzero coefficient vector,
and the same annihilation-operator range lemma works fiber by fiber. Orthogonal
complements then glue the fiberwise density into global density.

\section*{Proof}

We only spell out the points in \cite{LL} where finite dimensionality entered.
All references to equations and propositions in this proof are to the source
paper unless otherwise stated.

\subsection*{Boundedness and the kernel condition}

Since $K$ is bounded self-adjoint, the scale $\G_r$ is equivalent to the tensor
scale $\hhilb\otimes\FF_r$ for the values $r$ used below. Proposition 4.1 of
\cite{LL} gives $a(f_j)\in\mathcal B(\FF_s,\FF)$ and
$a^*(f_j)\in\mathcal B(\FF,\FF_{-s})$. Because there are only finitely many
bounded operators $B_j$, we have
\[
A\in\mathcal B(\G_s,\G),\qquad A^*\in\mathcal B(\G,\G_{-s}).
\]
The same coherent-vector proof as Proposition 4.3 of \cite{LL} gives density of
$\kerop(A|_{\G_2})$. Indeed, if
\[
\mathcal D_0=\{g\in\HH_2:\langle f_j,g\rangle=0\text{ for all }j\},
\]
then $\mathcal D_0$ is dense in $\HH$ by $\HH$-independence, and
$\hhilb\otimes\mathcal E(\mathcal D_0)$ is dense in $\G$ and contained in
$\kerop(A|_{\G_2})$.

\subsection*{Dense range via the joint spectral theorem}

Let $E$ be the joint spectral measure of the commuting normal tuple
$(B_1,\ldots,B_N)$ on its compact joint spectrum $\Sigma\subset\mathbb C^N$.
Choose a scalar spectral measure $\nu$ and a direct-integral representation
\[
\hhilb\simeq \int_\Sigma^\oplus \hhilb_\lambda\,\dd\nu(\lambda)
\]
in which $B_j$ acts as multiplication by the coordinate $\lambda_j$. The common
kernel hypothesis implies
\[
E(\{(0,\ldots,0)\})=0,
\]
because this projection is the kernel projection of
\(\sum_j B_j^*B_j\).

Tensoring with $\FF$ gives
\[
\G\simeq \int_\Sigma^\oplus (\hhilb_\lambda\otimes\FF)\,\dd\nu(\lambda).
\]
On the dense decomposable domain with Fock component in $\FF_2$, the operator
$A$ is decomposable. Its fiber at $\lambda$ is
\[
A(\lambda)=I_{\hhilb_\lambda}\otimes a(g_\lambda),
\]
where $g_\lambda$ is the corresponding nonzero linear combination of
$f_1,\ldots,f_N$ with coefficient vector $\lambda$ (up to the harmless conjugation
dictated by the inner-product convention). Since $\lambda\ne0$ for $\nu$-almost
every $\lambda$ and the family $f_1,\ldots,f_N$ is $\HH$-independent, we have
\[
g_\lambda\in \HH_{-s}\setminus\HH,\qquad g_\lambda\ne0
\]
for almost every $\lambda$.

By Lemma 4.2(ii) of \cite{LL}, $\ran(a(g_\lambda)|_{\FF_2})$ is dense in $\FF$
for almost every $\lambda$. Hence
\[
\ran\!\left(I_{\hhilb_\lambda}\otimes a(g_\lambda)
\bigm|_{\hhilb_\lambda\otimes\FF_2}\right)
\]
is dense in $\hhilb_\lambda\otimes\FF$ for almost every $\lambda$.

It remains only to pass from fibers to the direct integral. Let
$\Phi\in\G$ be orthogonal to $\ran(A|_{\G_2})$. Testing against decomposable
vectors supported on arbitrary measurable subsets of $\Sigma$ shows that, for
almost every $\lambda$, the fiber $\Phi(\lambda)$ is orthogonal to the dense
fiber range displayed above. Therefore $\Phi(\lambda)=0$ for almost every
$\lambda$, and $\Phi=0$. Thus $\ran(A|_{\G_2})^\perp=\{0\}$, i.e.
$\ran(A|_{\G_2})$ is dense in $\G$.

\subsection*{Application of Posilicano's theorem}

The two density requirements in Theorem 3.3/Theorem 3.13 of \cite{LL}, which
restate Posilicano's abstract results \cite{Pos}, are now verified. Since
$0<s\le1$, the operator $T=-ARA^*$ is bounded exactly as in Eq. (47) of
\cite{LL}. Therefore the proof of part (a) of Theorem 1.3 in \cite{LL} applies
verbatim and gives self-adjointness on the stated domain.

For norm resolvent convergence, choose $f_{j,n}\in\HH$ as in Proposition 4.1 of
\cite{LL}. The estimates in Eq. (74)--(76) of \cite{LL} use only the boundedness
of the finitely many $B_j$, not finite dimensionality. In the present setting,
for $\Psi\in\G_1$ and $A_n^\sharp=A_n$ or $A_n^\dagger$,
\[
\|A_n^\sharp\Psi\|
\le
\left(\sum_{j=1}^N \|B_j\|\,\|f_{j,n}\|_{\HH}\right)
\|I_{\hhilb}\otimes(\mathcal N+1)^{1/2}\Psi\|,
\]
and the lower bound $\omega\ge m_0>0$ again controls the last factor by the
$\G_1$ norm, up to a constant. Thus the regularized $A_n$ are closable and
infinitesimally $H_{\rm free}$-bounded, while
\[
\|A_n-A\|_{\G_s\to\G}
\le\sum_{j=1}^N \|B_j\|\,\|a(f_{j,n})-a(f_j)\|_{\FF_s\to\FF}\to0,
\]
and similarly for the adjoints. With $E_n=0$, the proof of part (b) of Theorem
1.3 in \cite{LL} applies unchanged, proving norm resolvent convergence.

\section*{Verification and Limitations}

The proof was checked against the source proof line by line. Finite
dimensionality is replaced only in Proposition 4.4, where the common eigenbasis
is replaced by the joint spectral theorem. The regularization estimates already
use operator norms $\|B_j\|$ and therefore survive unchanged for bounded
operators on a separable spin space.

This packet does not relax normality, pairwise commutation, or the zero common
kernel assumption. It also does not treat unbounded spin Hamiltonians $K$ or
the Case 2/Case 3 renormalization problems mentioned in the source.

Bounded novelty search on 2026-07-19 for exact phrases around
``generalized spin-boson'', ``mild ultraviolet'', ``infinite-dimensional'', and
``spectral calculus'' found the source paper and adjacent spin-boson literature,
but no later paper carrying out this precise bounded commuting-normal
infinite-dimensional extension.

\begin{thebibliography}{9}
\bibitem{LL}
Sascha Lill and Davide Lonigro,
\emph{Self-adjointness and domain of generalized spin-boson models with mild
ultraviolet divergences},
arXiv:2307.14727, 2023; version 2, 2025.
\url{https://arxiv.org/abs/2307.14727}

\bibitem{Pos}
Andrea Posilicano,
\emph{On the self-adjointness of $H+A^*+A$},
Mathematical Physics, Analysis and Geometry 23 (2020), no. 4, 1--31.
\end{thebibliography}

\end{document}
